\section{Related Work}\label{sec:related}

\begin{table*}[t]
  \centering
  \caption{Comparison of \sys{} with existing IoT simulation platforms, 3D
  simulators, and digital twin frameworks.  \ding{51} = supported,
  \ding{55} = not supported, \textbf{P} = partial.
  \sys{} is the only platform that integrates all six capabilities
  in a single reproducible framework.}
  \label{tab:system-comparison}
  \small
  \begin{tabular}{l c c c c c c}
    \toprule
    {Platform} & {3D Sim.} & {Firmware} & {Cloud IoT} & {Activity Gen.} & {Security} & {Open Source} \\
    \midrule
    IoTSim~\cite{iotsim}                 & \ding{55} & \ding{55} & \textbf{P}  & \ding{55} & \ding{55} & \ding{51} \\
    FIESTA-IoT~\cite{fiesta-iot}         & \ding{55} & \ding{55} & \ding{51}   & \ding{55} & \ding{55} & \ding{51} \\
    DPWSim~\cite{dpwsim}                 & \ding{55} & \ding{55} & \textbf{P}  & \ding{55} & \ding{55} & \ding{51} \\
    Habitat~3.0~\cite{habitat3}          & \ding{51} & \ding{55} & \ding{55}   & \ding{55} & \ding{55} & \ding{51} \\
    AI2-THOR~\cite{ai2thor}              & \ding{51} & \ding{55} & \ding{55}   & \ding{55} & \ding{55} & \ding{51} \\
    iGibson~2.0~\cite{igibson}           & \ding{51} & \ding{55} & \ding{55}   & \ding{55} & \ding{55} & \ding{51} \\
    Firmadyne~\cite{firmadyne}           & \ding{55} & \ding{51} & \ding{55}   & \ding{55} & \ding{51} & \ding{51} \\
    Azure DT~\cite{azure-digital-twins}  & \ding{55} & \ding{55} & \ding{51}   & \ding{55} & \ding{55} & \ding{55} \\
    AWS TwinMaker~\cite{aws-iot-twinmaker} & \ding{55} & \ding{55} & \ding{51} & \ding{55} & \ding{55} & \ding{55} \\
    \midrule
    \textbf{\sys{} (ours)} & \ding{51} & \ding{51} & \ding{51}   & \ding{51} & \ding{51} & \ding{51} \\
    \bottomrule
  \end{tabular}
  \vspace{-2mm}
\end{table*}


\Cref{tab:system-comparison} positions \sys{} against existing
platforms across six capability dimensions.  No prior system
integrates all six; \sys{} achieves this by composing Habitat~3.0
(3D simulation), Docker-based QEMU (firmware emulation), SmartThings
Schema Connector (cloud IoT), LLM scheduling (activity generation),
and a five-suite attack framework (security testing).

\paragraph{IoT Simulation Platforms.}
IoTSim~\cite{iotsim} models cloud-centric IoT workflows but lacks 3D
spatial grounding, firmware emulation, and occupant activity generation.
FIESTA-IoT~\cite{fiesta-iot} provides interoperability testing across
heterogeneous protocols but generates synthetic traffic rather
than behaviorally grounded events.
DPWSim~\cite{dpwsim} simulates DPWS device interactions at the
service layer without firmware-level fidelity.
SWoTSuite~\cite{swotsuite} benchmarks semantic web IoT applications
but does not model device firmware or physical environments.
In contrast, \sys{} unifies 3D embodied simulation, real firmware
execution, and cloud integration, enabling end-to-end testing from
human activity through firmware response to cloud state update.
Moreover, all attack traffic is captured as genuine \texttt{.pcap}
files (\Cref{sec:eval-pcap}), providing independently verifiable
evidence of network-realistic security testing---a level of
transparency not offered by any existing IoT simulator.

\paragraph{Smart Home Testbeds.}
The CASAS project~\cite{casas} maintains instrumented apartments
at Washington State University, providing longitudinal activity
datasets across multiple homes.  The Aware Home~\cite{aware-home}
at Georgia Tech and the PlaceLab~\cite{placelab} at MIT are seminal
physical testbeds.  These inspire \sys{}'s validation methodology
(we evaluate against CASAS and ARAS data) but are not reproducible
outside their host institutions, and they cannot be reconfigured to
test arbitrary floor plans, device configurations, or attack scenarios.

\paragraph{Activity Recognition Datasets.}
CASAS~\cite{casas}, ARAS~\cite{aras}, Opportunity~\cite{opportunity},
and SPHERE~\cite{sphere} provide real-world activity data.
\sys{} complements these by generating large-scale synthetic data
from diverse LLM-driven personas, achieving JS divergence of
0.101--0.113 against ARAS ground-truth labels (\Cref{sec:eval-rq1}).

\paragraph{Embodied AI Simulators.}
Habitat~\cite{habitat19iccv,habitat3} focuses on navigation and
object interaction for robotics.
AI2-THOR~\cite{ai2thor} and iGibson~\cite{igibson} offer similar
capabilities with interactive object state.
These platforms excel at visual perception and manipulation tasks
but lack IoT device modeling, firmware emulation, and cloud
connectivity---features central to smart home research.
\sys{} extends Habitat~3.0 with these capabilities while preserving
its physics-accurate humanoid navigation.

\paragraph{LLM-Based Agents.}
Park et al.~\cite{generative-agents} demonstrated LLM-driven agents
in a virtual town, showing that language models can produce believable
daily routines.  Voyager~\cite{voyager} uses LLMs for open-ended
exploration in Minecraft.  Both operate in unconstrained sandbox
environments.  \sys{} applies LLM-driven behavior generation to the
\emph{structured} smart home domain, with quantitative validation
against real-world activity datasets (CASAS, ARAS) and a formal
baseline comparison against Markov-chain activity generators.

\paragraph{Firmware Analysis.}
Firmadyne~\cite{firmadyne} and Avatar~\cite{avatar} use QEMU for
automated firmware analysis and vulnerability discovery in
Linux-based IoT firmware.
\sys{} repurposes QEMU emulation for \emph{behavioral} simulation
rather than standalone security analysis, coupling firmware execution
with 3D occupant simulation and cloud synchronization to enable
integrated firmware-in-the-loop testing within realistic usage
scenarios.

\paragraph{IoT Protocol Attacks.}
Fu~et~al.~\cite{phantom-delay} demonstrated phantom-delay attacks
that exploit IoT timeout behaviors to stealthily delay device events
for up to 40~minutes on Apple HomeKit and 64~seconds on Ring.
Their work required physical IoT deployments and manual network
configuration.  \sys{} enables reproduction of such protocol-level
attacks within a simulated environment
(\Cref{sec:eval-phantom-delay}), reducing the barrier to entry for
IoT security research.  Our CVE validation (\Cref{tab:cve-validation})
further demonstrates that the platform's 36~attack implementations
cover vulnerability classes representative of 24~real-world CVEs.

\paragraph{Digital Twins for IoT.}
Azure Digital Twins~\cite{azure-digital-twins} and AWS IoT
TwinMaker~\cite{aws-iot-twinmaker} provide cloud-based digital twin
frameworks for IoT data modeling and visualization.  These platforms
enable device state mirroring and analytics dashboards but lack
behavioral simulation with embodied agents, firmware-in-the-loop
execution, and integrated security testing---all of which \sys{}
provides as a self-contained open-source platform.
